\section{Introduction}
\label{sec:intro}

Distinguishing between texts composed by human writers and those generated by large language models
(LLMs) has emerged as one of the most consequential challenges in natural language processing.
The rapid proliferation of capable autoregressive models—ranging from GPT-4 to Falcon and Llama
variants—has dramatically lowered the cost of synthesising fluent, topically coherent prose, which
in turn threatens the integrity of academic assessment, online discourse, and information ecosystems
at large~\cite{jakesch:2023,ippolito:2020}.

The \textit{Voight-Kampff Generative AI Detection} task at PAN 2026 encapsulates precisely this
challenge in a rigorous shared-task setting~\cite{bevendorff:2025}.  Participants receive a
single text—potentially subjected to adversarial obfuscation—and must return a confidence score
in $[0,1]$ indicating the probability that it was machine-authored.  Crucially, the task introduces
a \emph{builder–breaker} dynamic: while PAN participants construct detectors, ELOQUENT participants
actively design text-generation and obfuscation strategies intended to defeat them.  This adversarial
framing substantially raises the bar beyond conventional machine-generated text (MGT) detection
benchmarks.

Two complementary paradigms have emerged for LLM detection~\cite{bevendorff:2025b}.  The first
treats detection as an \emph{authorship attribution} problem: given a fixed reference corpus, the
system must decide which author—human or machine—produced the query text.
The second, adopted here, frames it as an \emph{authorship verification} problem, querying only
whether the text was human-authored, without any reference sample.  Verification is inherently
harder because detectors must generalise across writing styles, topics, and adversarial perturbations
without author-specific anchors.

Zero-shot detectors such as DetectGPT~\cite{mitchell:2023}, Binoculars~\cite{hans:2024}, and
DetectLLM~\cite{su:2023} exploit the statistical properties of LM output distributions but
degrade noticeably when texts are paraphrased or subjected to character-level
obfuscation~\cite{sadasivan:2023}.  Supervised classifiers trained on TF-IDF features achieve
strong results on in-domain data~\cite{bevendorff:2024d} yet remain brittle under distribution
shift.  A meaningful advance requires representations that are simultaneously rich enough to
distinguish subtle stylistic signatures and robust enough to withstand deliberate interference.

In this paper we introduce \radar{} (\textbf{R}obust \textbf{A}dversarial-Resistant \textbf{D}etection with
\textbf{A}daptive \textbf{R}easoning), a multi-modal fusion architecture for AI text detection submitted to
PAN 2026.  \radar{} combines:
\begin{enumerate}
  \item \textbf{\arse{}} (\textit{Adaptive Robust Semantic Encoder}): a \modernbert{} transformer
        encoder that captures deep contextual and semantic representations of the input text.
  \item \textbf{\slpe{}} (\textit{Statistical Linguistic Profile Extractor}): a hand-crafted,
        computation-efficient module that computes 38 interpretable stylometric features
        organised across five linguistic categories.
  \item \textbf{Cross-Attention Fusion}: a multi-head cross-attention layer that dynamically
        integrates semantic and stylometric signals.
  \item \textbf{\uacc{}} (\textit{Unified Adaptive Classification Controller}): a single linear
        head that maps the fused representation to a calibrated detection score.
\end{enumerate}

By pairing deep semantic encoding with explicit stylometric evidence, \radar{} addresses the core
weakness of purely neural approaches: susceptibility to adversarial style mimicry.  Our two-stage
training protocol—frozen-encoder pre-training followed by full end-to-end fine-tuning—enables
stable optimisation of the 400M-parameter model on a moderately sized combined corpus.

Experiments on the official PAN-2026 validation set and the external RAID benchmark demonstrate
that \radar{} achieves a mean score of \textbf{0.978} across all five official evaluation metrics,
representing a substantial improvement over the TF-IDF SVM baseline (mean $= 0.746$ on the same
combined evaluation set) and competitive parity with the official PAN-2026 TF-IDF baseline
(mean $= 0.978$) which, however, does not generalise under distribution shift.

The remainder of the paper is organised as follows.
Section~\ref{sec:related} reviews related work.
Section~\ref{sec:method} describes the \radar{} architecture in detail.
Section~\ref{sec:data} presents the datasets and preprocessing pipeline.
Section~\ref{sec:experiments} describes the experimental setup and training procedure.
Section~\ref{sec:results} reports and analyses the results.
Section~\ref{sec:conclusion} concludes.
